\documentclass[11pt,a4paper]{article}

\usepackage[spanish,es-noshorthands]{babel}
\usepackage[utf8]{inputenc}
\usepackage[T1]{fontenc}
\usepackage{hyperref}
\usepackage{geometry}
\usepackage{longtable}
\usepackage{graphicx}
\usepackage{listings}
\usepackage{tikz}
\usetikzlibrary{positioning}


\geometry{margin=2.5cm}
\hypersetup{
    colorlinks=true,
    linkcolor=blue,
    urlcolor=blue
}

\title{README -- Proyecto Web Quixo}
\author{GRUPO 4 \\ Programación Avanzada Web}
\date{\today}

\begin{document}

\maketitle

%====================================================
\section{Información general del proyecto}
%====================================================

\begin{description}
    \item[Número de grupo:] \textbf{4}
    \item[Nombre del curso:] \textbf{Programación Avanzada Web}
    \item[Nombre de la aplicación:] \textbf{Proyecto Web Quixo}
\end{description}

Breve descripción del proyecto:

\medskip

La aplicación web desarrollada implementa el juego de mesa \textbf{Quixo}, una evolución del clásico Tic-tac-toe (``gato''). 
Permite jugar en modo \textbf{dos jugadores (2P)} y \textbf{cuatro jugadores (4P)}, visualizar y exportar partidas finalizadas,
así como consultar estadísticas de jugadores y equipos.

%====================================================
\section{Integrantes del grupo}
%====================================================

A continuación se muestran los integrantes del grupo con nombre completo y carné:

\begin{longtable}{|l|l|}
\hline
\textbf{Nombre completo}      & \textbf{Carné} \\
\hline
Sebastián Mora Figueroa        & FI22027357 \\ \hline
Kevin Sánchez Seas             & FI22025309 \\ \hline
Fiorella Portuguez Rojas       & FI22024660 \\ \hline
\end{longtable}

%====================================================
\section{Información de Git de los integrantes}
%====================================================

En la siguiente tabla se detalla la información de Git de cada integrante 
(usuario y correo de Git):

\begin{longtable}{|l|l|l|}
\hline
\textbf{Nombre}                 & \textbf{Usuario Git} & \textbf{Correo Git}              \\
\hline
Sebastián Mora Figueroa        & samorafi             & sebasmora172522@gmail.com        \\ \hline
Kevin Sánchez Seas             & kasanche11           & ksanchez90688@ufide.ac.cr        \\ \hline
Fiorella Portuguez Rojas       & fioportu             & portuguezfiore@gmail.com         \\ \hline
\end{longtable}


%====================================================
\section{Tecnologías, tipo de aplicación y arquitectura}
%====================================================

\subsection{Frameworks y herramientas utilizadas}

A continuación se listan los frameworks y herramientas principales utilizados:

\begin{itemize}
    \item \textbf{Tipo de aplicación:} SPA (Single Page Application).
    \item \textbf{Frontend:} React (Vite).
    \item \textbf{Backend:} ASP.NET Core (Web API).
    \item \textbf{ORM:} Entity Framework Core.
    \item \textbf{Base de datos:} SQL Server.
    \item \textbf{Cliente HTTP:} Axios (frontend) para consumir la API.
    \item \textbf{Enrutamiento en frontend:} React Router DOM.
    \item \textbf{Editor/IDE:} Visual Studio Code
    \item \textbf{Control de versiones:} GitHub.
\end{itemize}


\subsection{Tipo de aplicación: SPA}

La aplicación es una \textbf{SPA (Single Page Application)} porque:
\begin{itemize}
    \item El frontend en React se carga una sola vez en el navegador.
    \item La navegación entre vistas (menú principal, juego, historial, estadísticas) se realiza de forma dinámica en el cliente,
          usando React Router, sin recargar completamente la página.
    \item La comunicación con el servidor se realiza mediante peticiones a una Web API (REST) para obtener y enviar datos (jugadores, partidas, estadísticas).
\end{itemize}

\subsection{Arquitectura utilizada: Web API + SPA}

La arquitectura general combina:
\begin{itemize}
    \item \textbf{Backend ASP.NET Core Web API} con una estructura por capas (controladores, capa de dominio/servicios, acceso a datos con EF Core).
    \item \textbf{Frontend SPA en React} que consume los endpoints de la Web API.
    \item \textbf{Base de datos SQL Server} para persistir jugadores, partidas, historial y estadísticas.
\end{itemize}

Razones para utilizar esta arquitectura:
\begin{itemize}
    \item Separación clara de responsabilidades entre backend (lógica de negocio, persistencia) y frontend (interfaz de usuario).
    \item Facilidad para exponer servicios REST que podrían ser consumidos por otras interfaces en el futuro.
    \item React y ASP.NET Core son tecnologías ampliamente utilizadas, con buena documentación y soporte.
\end{itemize}

%====================================================
\section{Definición de la base de datos}
%====================================================

La aplicación utiliza una base de datos relacional en \textbf{SQL Server}. 
A continuación se muestra el script simplificado de las tablas principales:

\subsection{Script de creación de tablas}

\begin{lstlisting}[language=SQL]
-- TABLA: Jugadores
CREATE TABLE dbo.Jugadores (
  JugadorId       INT IDENTITY(1,1) NOT NULL PRIMARY KEY,
  Nombre          NVARCHAR(100)     NOT NULL,
  FechaCreacion   DATETIME2(3)      NOT NULL DEFAULT SYSUTCDATETIME()
);

-- TABLA: Partidas
CREATE TABLE dbo.Partidas (
  PartidaId         INT IDENTITY(1,1) NOT NULL PRIMARY KEY,
  Modo              CHAR(2)           NOT NULL
    CONSTRAINT CK_Partidas_Modo CHECK (Modo IN ('2P','4P')),
  FechaCreacion     DATETIME2(0)      NOT NULL DEFAULT SYSDATETIME(),
  FechaSobrescrita  DATETIME2(0)      NULL,
  FechaFinalizada   DATETIME2(0)      NULL,
  DuracionSegundos  INT               NULL,
  JugadorOid        INT               NULL,
  JugadorXid        INT               NULL,
  EquipoA1Id        INT               NULL,
  EquipoA2Id        INT               NULL,
  EquipoB1Id        INT               NULL,
  EquipoB2Id        INT               NULL,
  GanadorSimbolo    CHAR(1)           NULL,  -- 'O' o 'X' (solo 2P)
  GanadorEquipo     CHAR(1)           NULL,  -- 'A' o 'B' (solo 4P)
  TableroFinalXml   XML               NULL,
  HistorialXml      XML               NOT NULL
);

-- TABLA: Equipos4P
CREATE TABLE dbo.Equipos4P (
  EquipoId       INT IDENTITY(1,1) NOT NULL PRIMARY KEY,
  Jugador1Id     INT               NOT NULL,
  Jugador2Id     INT               NOT NULL,
  FechaCreacion  DATETIME2(3)      NOT NULL DEFAULT SYSUTCDATETIME()
);

-- TABLA: PartidasEquipos4P
CREATE TABLE dbo.PartidasEquipos4P (
  PartidaEquipoId INT IDENTITY(1,1) NOT NULL PRIMARY KEY,
  PartidaId       INT               NOT NULL,   -- Modo = '4P'
  EquipoId        INT               NOT NULL,   -- Referencia a Equipos4P
  Letra           CHAR(1)           NOT NULL
    CONSTRAINT CK_PartidasEquipos4P_Letra CHECK (Letra IN ('A','B')),
  Simbolo         CHAR(1)           NOT NULL
    CONSTRAINT CK_PartidasEquipos4P_Simbolo CHECK (Simbolo IN ('O','X'))
);

-- TABLA: PartidaParticipantes
CREATE TABLE dbo.PartidaParticipantes (
  PartidaParticipanteId INT IDENTITY(1,1) NOT NULL PRIMARY KEY,
  PartidaId             INT               NOT NULL,
  JugadorId             INT               NOT NULL,
  Rol                   NVARCHAR(3)       NOT NULL,
  Equipo                CHAR(1)           NULL
    CONSTRAINT CK_PartidaParticipantes_Equipo 
        CHECK (Equipo IN ('A','B') OR Equipo IS NULL),
  OrdenTurno            SMALLINT          NULL
);

-- TABLA: EstadisticasJugadores2P
CREATE TABLE dbo.EstadisticasJugadores2P (
  JugadorId       INT           NOT NULL PRIMARY KEY,  
  Jugadas         INT           NOT NULL DEFAULT 0,
  Ganadas         INT           NOT NULL DEFAULT 0,
  EfectividadPorc DECIMAL(5,2)  NULL
);

-- TABLA: EstadisticasEquipos4P
CREATE TABLE dbo.EstadisticasEquipos4P (
  EquipoId        INT           NOT NULL PRIMARY KEY,  
  Jugadas         INT           NOT NULL DEFAULT 0,
  Ganadas         INT           NOT NULL DEFAULT 0,
  EfectividadPorc DECIMAL(5,2)  NULL
);
\end{lstlisting}

\subsection{Descripción de entidades}

\begin{itemize}
    \item \textbf{Jugadores:} almacena los datos básicos de cada jugador (identificador y nombre).
    \item \textbf{Partidas:} guarda la información general de cada partida (modo 2P o 4P, fechas, duración, 
          participantes, ganador, tablero final e historial completo en XML).
    \item \textbf{Equipos4P:} define las parejas reutilizables para el modo de cuatro jugadores.
    \item \textbf{PartidasEquipos4P:} relaciona qué equipos (A y B) jugaron una partida 4P y qué símbolo usaron (O o X).
    \item \textbf{PartidaParticipantes:} lista detallada de los participantes en cada partida, 
          unificando los roles:
          \begin{itemize}
              \item Modo 2P: ``O'' y ``X''.
              \item Modo 4P: ``A1'', ``A2'', ``B1'', ``B2'', con su respectivo equipo A o B y orden de turno.
          \end{itemize}
    \item \textbf{EstadisticasJugadores2P:} estadísticas agregadas para modo de 2 jugadores (jugadas, ganadas, efectividad).
    \item \textbf{EstadisticasEquipos4P:} estadísticas agregadas para equipos en modo 4 jugadores.
\end{itemize}

\subsection{Diagrama de base de datos}

El diagrama completo de la base de datos puede consultarse en el siguiente enlace:

\begin{itemize}
    \item \url{https://github.com/samorafi/Proyecto_PrograAvanzadaWeb/blob/main/DB/DBQuixo.png}
\end{itemize}

%====================================================
\section{Instructivo de instalación, compilación y ejecución}
%====================================================

\subsection{Backend: ASP.NET Core Web API}

\subsubsection{Requisitos previos}

\begin{itemize}
    \item .NET SDK (versión utilizada en el proyecto .NET 8.0).
    \item SQL Server instalado y con una base de datos accesible para \texttt{QuixoDb}.
    \item Herramienta para ejecutar scripts SQL (SQL Server Management Studio).
\end{itemize}

\subsubsection{Crear solución y proyecto}

\begin{lstlisting}[language=bash]
# Carpeta de trabajo
mkdir quixo && cd quixo

# Solucion y proyecto Web API
dotnet new sln -n Quixo
dotnet new webapi -n Quixo.Api

# Agregar proyecto a la solucion
dotnet sln Quixo.sln add Quixo.Api/Quixo.Api.csproj
\end{lstlisting}

\subsubsection{Instalación de paquetes NuGet}

\begin{lstlisting}[language=bash]
cd Quixo.Api

# EF Core para SQL Server y herramientas de migraciones
dotnet add package Microsoft.EntityFrameworkCore.SqlServer
dotnet add package Microsoft.EntityFrameworkCore.Tools
dotnet add package Microsoft.EntityFrameworkCore
\end{lstlisting}


\subsubsection{Configuración de la base de datos}

\begin{enumerate}
    \item Crear la base de datos \texttt{QuixoDb} en SQL Server.
    \item Ejecutar el script de creación de tablas mostrado en la Sección~4.
    \item Configurar la cadena de conexión en \texttt{appsettings.json}:
\end{enumerate}

\begin{lstlisting}
"ConnectionStrings": {
  "SqlServer": "Server=SERVIDOR;Database=QuixoDb;Trusted_Connection=True;TrustServerCertificate=True"
}
\end{lstlisting}


\subsubsection{Compilación y ejecución del backend}

\begin{lstlisting}[language=bash]
# Desde la carpeta Quixo.Api
dotnet build
dotnet run
\end{lstlisting}

El backend levantará la Web API (\texttt{http://localhost:5207}).

\subsection{Frontend: React (SPA con Vite)}

\subsubsection{Creación e instalación inicial}

\begin{lstlisting}[language=bash]
npm create vite@latest quixo-front -- --template react
# Framework: React
# Variant: JavaScript
# rolldown-vite: No

cd quixo-front
npm install
npm i axios react-router-dom classnames
\end{lstlisting}

\subsubsection{Configuración}

\begin{itemize}
    \item Configurar la URL base de la API (\texttt{http://localhost:5207}) en el archivo de servicios de Axios.
    \item Ajustar rutas de React Router para:
    \begin{itemize}
        \item Menú principal (jugar nueva partida, ver partidas finalizadas, ver estadísticas).
        \item Vista de juego (tablero 5x5, reloj de 7 segmentos, botón de reinicio).
        \item Vista de historial de partidas.
        \item Vista de estadísticas de jugadores y equipos.
    \end{itemize}
\end{itemize}

\subsubsection{Compilación y ejecución del frontend}

\begin{lstlisting}[language=bash]
# Desarrollo
npm run dev

# Compilacion para produccion
npm run build
\end{lstlisting}

La aplicación esta disponible en \texttt{http://localhost:5173}.

%====================================================
\section{Referencias (sitios web y material consultado)}
%====================================================

A continuación se listan las referencias de sitios web utilizados para entender reglas, conceptos y ejemplos:

\begin{itemize}
    \item \textbf{Cómo jugar Quixo:}\\
    \url{https://maldon.com.ar/blog/projects/quixo/}

    \item \textbf{QUIXO - ¿Cómo se juega? - Tutorial en español y partida:}\\
    \url{https://www.youtube.com/watch?v=IHtfOYiNuEo}
    
    \item \textbf{SPA vs MPA y las arquitecturas Web:}\\
    \url{https://www.arquitecturajava.com/spa-vs-mpa-y-las-arquitecturas-web/}
    
    \item \textbf{Tutorial: Creación de aplicaciones ASP.NET Core con React:}\\
    \url{https://learn.microsoft.com/es-es/visualstudio/javascript/tutorial-asp-net-core-with-react?view=visualstudio}
    
    \item \textbf{Display de 7 segmentos (conceptos y funcionamiento):}\\
    \url{https://industriasgsl.com/blogs/automatizacion/display-de-7-segmentos}

    \item \textbf{Language Integrated Query (LINQ):}\\
    \url{https://learn.microsoft.com/es-mx/dotnet/csharp/linq/}

    \item \textbf{Habilitar solicitudes entre orígenes (CORS) en ASP.NET Core:}\\
    \url{https://learn.microsoft.com/es-mx/aspnet/core/security/cors?view=aspnetcore-10.0}

    \item \textbf{Fetch API y llamadas al backend:}\\
    \url{https://developer.mozilla.org/es/docs/Web/API/Fetch_API/Using_Fetch}

    \item \textbf{Timers en React y técnicas recomendadas:}\\
    \url{https://react.dev/learn/keeping-components-pure}

    \item \textbf{Crear un documento en LaTeX:}\\
    \url{https://manualdelatex.com/tutoriales/crear-un-documento}
\end{itemize}

%====================================================
\section{Uso de agentes de IA y prompts}
%====================================================

En el desarrollo del proyecto se utilizaron agentes de IA para apoyo en el diseño de componentes específicos.
A continuación se listan los prompts y recursos principales:

\subsection{Prompts y enlaces}

\begin{itemize}
    \item \textbf{Cómo hacer un tablero 5x5 con React (X, O y neutro):}\\
    Prompt y conversación: \url{https://chat.deepseek.com/share/0uqkmm2gqnp4z6zftn}
    
    \item \textbf{Reloj digital de 7 segmentos en React:}\\
    Prompt y conversación: \url{https://chat.deepseek.com/share/0uqkmm2gqnp4z6zftn}
    
    \item \textbf{Fórmula de efectividad (porcentaje de partidas ganadas):}\\
    Prompt y conversación: \url{https://chat.deepseek.com/share/0uqkmm2gqnp4z6zftn}

    \item \textbf{Generar tablas con migraciones}\\
    Prompt y conversación: \url{https://chatgpt.com/share/69373796-790c-8002-93a5-3536469bd2bd}

    \item \textbf{Modificar lógica para 2P y 4P}\\
    Prompt y conversación: \url{https://chatgpt.com/share/69373ac8-964c-8002-8fe2-3491c0959d94}
\end{itemize}

%====================================================
\end{document}
